% !TEX program = xelatex
%%%%%%%%%%%%%%%%%%%%%%%%%%%%%%%%%%%%%%%%%%%%%%%%%%%%%%%%%%%%%%%%%%%%%%%%%%%%%%%
%  SCX-Audited Medical Diagnosis:
%  Multi-Physician Consensus for Certified Clinical Decision-Making
%  SCX审计医学诊断：多医师共识认证的临床决策
%%%%%%%%%%%%%%%%%%%%%%%%%%%%%%%%%%%%%%%%%%%%%%%%%%%%%%%%%%%%%%%%%%%%%%%%%%%%%%%
%  Chinese terminology:
%    医学诊断 (Medical Diagnosis) / 多学科会诊 (Multidisciplinary Consultation)
%    误诊率 (Misdiagnosis Rate) / 第二意见 (Second Opinion)
%    肿瘤委员会 (Tumor Board) / 影像学 (Imaging) / 病理学 (Pathology)
%    检验医学 (Laboratory Medicine) / 鉴别诊断 (Differential Diagnosis)
%%%%%%%%%%%%%%%%%%%%%%%%%%%%%%%%%%%%%%%%%%%%%%%%%%%%%%%%%%%%%%%%%%%%%%%%%%%%%%%
\documentclass[12pt,a4paper]{ctexart}

% ── Packages ────────────────────────────────────────────────────────────────
\usepackage{amsmath,amssymb,amsthm}
\usepackage{mathtools}
\usepackage{bm}
\usepackage{booktabs}
\usepackage{graphicx}
\usepackage{xcolor}
\usepackage{hyperref}
\usepackage{enumitem}
\usepackage[margin=2.5cm]{geometry}
\usepackage{setspace}
\usepackage{dsfont}
\usepackage{multirow}
\usepackage{array}
\usepackage{cleveref}
\usepackage{algorithm}
\usepackage{algpseudocode}
\usepackage{bbm}
\usepackage{longtable}
\onehalfspacing

% ── Theorem Environments ────────────────────────────────────────────────────
\newtheorem{theorem}{Theorem}[section]
\newtheorem{lemma}[theorem]{Lemma}
\newtheorem{corollary}[theorem]{Corollary}
\newtheorem{proposition}[theorem]{Proposition}
\newtheorem{definition}[theorem]{Definition}
\newtheorem{remark}[theorem]{Remark}
\theoremstyle{definition}
\newtheorem{assumption}[theorem]{Assumption}
\newtheorem{protocol}[theorem]{Protocol}

% ── Rigor Labels ────────────────────────────────────────────────────────────
\definecolor{rigorFull}{RGB}{0,100,0}
\definecolor{rigorPartial}{RGB}{180,120,0}
\definecolor{rigorSketch}{RGB}{150,0,0}
\newcommand{\rigorFull}{{\color{rigorFull}$\blacksquare$\,\textbf{[Rigor: FULL $\star\star\star$]}}}
\newcommand{\rigorPartial}{{\color{rigorPartial}$\blacksquare$\,\textbf{[Rigor: PARTIAL $\star\star\;$]}}}
\newcommand{\rigorSketch}{{\color{rigorSketch}$\blacksquare$\,\textbf{[Rigor: SKETCH $\star\;\,$]}}}

% ── SCX Framework Macros ────────────────────────────────────────────────────
\newcommand{\Situs}{\textsc{Situs}}
\newcommand{\Yajie}{\textsc{Yajie}}
\newcommand{\Spring}{\textsc{Spring}}
\newcommand{\Cercis}{\textsc{Cercis}}
\newcommand{\SCX}{\textsc{SCX}}

% ── Math Operators ──────────────────────────────────────────────────────────
\DeclareMathOperator*{\argmax}{arg\,max}
\DeclareMathOperator*{\argmin}{arg\,min}
\DeclareMathOperator{\KL}{KL}
\DeclareMathOperator{\TV}{TV}
\DeclareMathOperator{\Var}{Var}
\DeclareMathOperator{\Cov}{Cov}
\DeclareMathOperator{\Bias}{Bias}
\DeclareMathOperator{\MSE}{MSE}
\DeclareMathOperator{\RMSE}{RMSE}
\DeclareMathOperator{\sign}{sign}
\DeclareMathOperator{\diag}{diag}
\DeclareMathOperator{\supp}{supp}
\DeclareMathOperator{\Lip}{Lip}
\DeclareMathOperator{\AUC}{AUC}
\DeclareMathOperator{\Sens}{Sens}
\DeclareMathOperator{\Spec}{Spec}
\DeclareMathOperator{\PPV}{PPV}
\DeclareMathOperator{\NPV}{NPV}

% ── Notation ────────────────────────────────────────────────────────────────
\newcommand{\R}{\mathbb{R}}
\newcommand{\N}{\mathbb{N}}
\newcommand{\E}{\mathbb{E}}
\newcommand{\Pbb}{\mathbb{P}}
\newcommand{\I}{\mathbb{I}}
\newcommand{\ind}[1]{\mathbf{1}[#1]}
\newcommand{\norm}[1]{\left\lVert#1\right\rVert}
\newcommand{\abs}[1]{\left|#1\right|}
\newcommand{\inner}[2]{\langle #1, #2 \rangle}
\newcommand{\cD}{\mathcal{D}}
\newcommand{\cF}{\mathcal{F}}
\newcommand{\cM}{\mathcal{M}}
\newcommand{\cP}{\mathcal{P}}
\newcommand{\cS}{\mathcal{S}}
\newcommand{\cX}{\mathcal{X}}
\newcommand{\cY}{\mathcal{Y}}
\newcommand{\cZ}{\mathcal{Z}}

% ── Medical Macros ──────────────────────────────────────────────────────────
\newcommand{\MisDiag}{误诊率}
\newcommand{\SecondOpinion}{第二意见}
\newcommand{\TumorBoard}{肿瘤委员会}
\newcommand{\MultiDiscConsult}{多学科会诊}
\newcommand{\MedDiag}{医学诊断}
\newcommand{\DiffDiag}{鉴别诊断}
\newcommand{\ImagingMod}{影像学}
\newcommand{\PathologyMod}{病理学}
\newcommand{\LabMod}{检验医学}

% ── Assumption Tags ─────────────────────────────────────────────────────────
\newcommand{\assumptionTag}[1]{\textbf{[#1]}}

% ── Title ───────────────────────────────────────────────────────────────────
\title{\textbf{SCX-Audited Medical Diagnosis:}\\[4pt]
       \large Multi-Physician Consensus for Certified Clinical Decision-Making\\[8pt]
       \normalsize \textbf{SCX审计\MedDiag{}：多医师共识认证的临床决策}}
\author{SCX Medical Working Group\\[4pt]
        \small Xiaogan Supercomputing Center \& Nous Research}
\date{June 2026}

\begin{document}

\maketitle

% ===========================================================================
\begin{abstract}
\noindent
Medical diagnosis (\MedDiag) is a high-stakes prediction task where error — misdiagnosis (\MisDiag) — carries direct human cost. Yet clinical practice provides no formal guarantee that diagnostic errors will be detected before harm occurs. We formalize medical diagnosis within the \SCX{} audit framework, treating physicians as expert agents ($f_m \in \cF$) making diagnostic claims ($\gamma_i(s)$) over a shared patient state $s$, and the diagnostic process as a multi-expert consensus problem.

Our central result is a \textbf{multi-physician error detection bound} (Theorem~\ref{thm:multi_diag}): with $M$ physicians of effective multiplicity $M_{\text{eff}}$, the probability that all $M$ miss a diagnostic error of clinical magnitude $\Delta$ satisfies $\Pbb(\text{all miss} \mid \text{error}) \leq \exp(-2M_{\text{eff}}\Delta^2)$. This formalizes the commonplace intuition of the \SecondOpinion{} ($M=2$): a second physician reduces misdiagnosis probability from $p$ to $p^2 / (1 + \bar{\rho})$, where $\bar{\rho}$ is inter-physician error correlation. We extend this to the \TumorBoard{} (\MultiDiscConsult), which instantiates the \Yajie{} multi-expert consensus protocol with $M \geq 5$ specialists from distinct modalities — imaging (\ImagingMod), pathology (\PathologyMod), laboratory medicine (\LabMod), surgery, and oncology — each contributing an independently-trained diagnostic opinion.

We prove three additional theorems: \textbf{Theorem~\ref{thm:self_diag}} (Self-Diagnosis Prohibition, 自诊断禁止定理) establishes that a physician auditing their own diagnosis provides zero epistemic information; \textbf{Theorem~\ref{thm:unident}} (Diagnostic Unidentifiability, 诊断不可区分定理) proves that when physicians disagree, the source — misdiagnosis vs.\ atypical disease presentation — is fundamentally unidentifiable without declared domain assumptions; and \textbf{Theorem~\ref{thm:situs_med}} (\Situs{} Multi-Modal Encoding, \Situs{}多模态编码定理) certifies the quality of fused imaging, laboratory, and pathology encodings via a Lipschitz-bounded representation space.

We introduce the \Cercis{} Score for diagnostic quality $S = Q + \eta N$, where $Q$ aggregates consensus reliability (inter-rater agreement, autopsy confirmation rate, clinical outcome concordance) and $N$ quantifies diagnostic novelty (atypical presentations, rare diseases, novel clinical patterns). The \Spring{} module permanently records every diagnostic decision, expert opinion, and outcome, enabling retrospective audit and cumulative improvement. We provide a concrete experimental protocol benchmarking against MIMIC-IV, CheXpert, and institutional tumor board records. Our framework yields a falsifiable prediction: hospitals implementing \SCX{}-audited diagnosis with $M \geq 5$ will exhibit $\MisDiag{}$ rates at least $1 - \exp(-2M_{\text{eff}}\Delta^2)$ lower than those with solo diagnosis, controlling for case complexity.

\medskip
\noindent\textbf{医学诊断是高风险的预测任务——误诊直接导致患者伤害。本文将医学诊断形式化为SCX审计框架下的多专家共识问题：医师即专家，诊断即声明，第二意见即M=2审计，肿瘤多学科会诊即Yajie共识。我们证明：M位医师的误诊漏检概率有Hoeffding指数界，自诊断等价于无审计，医师分歧时误诊与非典型表现不可区分（除非声明领域假设），并提出Situs多模态编码和Cercis诊断质量评分。实验方案覆盖MIMIC-IV、CheXpert和肿瘤委员会记录。}

\medskip
\noindent\textbf{Keywords:} medical diagnosis, multi-physician consensus, misdiagnosis rate, second opinion, tumor board, SCX audit, clinical decision-making, Cercis Score, \MedDiag, \MultiDiscConsult, \MisDiag, \SecondOpinion, \TumorBoard
\end{abstract}

\tableofcontents
\newpage

% ===========================================================================
\section{Introduction 引言}
\label{sec:intro}
% ===========================================================================

Medical error is the third leading cause of death in the United States, with diagnostic error — misdiagnosis (\MisDiag) — accounting for an estimated 40,000--80,000 deaths annually \cite{makary2016medical, graber2013incidence}. In China, studies report \MisDiag{} rates of 27\%--42\% in outpatient settings and 20\%--30\% in inpatient care \cite{wang2019misdiagnosis}. Globally, the World Health Organization estimates that diagnostic errors affect 5\% of adults in outpatient care annually. The problem is not one of physician competence — it is a structural absence of formal audit in the diagnostic process.

Consider the standard clinical workflow: a patient presents with symptoms $x$; a physician $f_1$ renders diagnosis $\hat{y}_1 = f_1(x)$. If the diagnosis is correct, treatment proceeds. If incorrect, the error may be discovered only after treatment failure, disease progression, or autopsy — all costly, some irreversible. The \SecondOpinion{} (\SecondOpinion) exists as an informal corrective: consult physician $f_2$, compare diagnoses. But the second opinion is typically requested only when doubt arises — and doubt is itself a diagnostic judgment made by the same system whose error it seeks to detect.

The \TumorBoard{} (\MultiDiscConsult) is the most mature instantiation of multi-expert consensus in clinical medicine. Tumor boards convene radiologists, pathologists, surgeons, medical oncologists, and radiation oncologists — typically $M = 5$--$12$ specialists — to reach a consensus diagnosis and treatment plan. Studies show tumor boards change the diagnosis in 25\%--52\% of cases and alter management in 30\%--65\% \cite{pillay2016impact, lamb2011quality}. Yet tumor boards are ad-hoc: they lack formal consensus protocols, make no quantitative error guarantees, and operate without permanent auditable records of dissenting opinions.

The \SCX{} framework provides the missing formalism. We map the medical diagnostic process onto the SCX multi-expert audit architecture:

\begin{center}
\begin{tabular}{l l}
\toprule
\textbf{Medical Concept} & \textbf{SCX Formalization} \\
\midrule
Physician & Expert $f_m \in \cF = \{f_1, \ldots, f_M\}$ \\
Diagnosis & Claim $\gamma_m(s) = (d_m(s), \sigma_m(s))$ \\
Patient state & $s \in \cS$ (history, exam, labs, imaging, pathology) \\
Second Opinion & $M=2$ audit with effective multiplicity $M_{\text{eff}}$ \\
Tumor Board & \Yajie{} multi-expert consensus with $M \geq 5$ \\
Misdiagnosis rate & $\Pbb(\text{all } M \text{ miss error} \mid \text{error})$ \\
Imaging/Lab/Pathology & Multi-modal \Situs{} encoding \\
Autopsy confirmation & Ground truth $g(s)$ (post-hoc verification) \\
Atypical presentation & Distribution shift $\cD_{\text{train}} \to \cD_{\text{novel}}$ \\
\bottomrule
\end{tabular}
\end{center}

\subsection{The Diagnostic Audit Gap 诊断审计空白}

Despite decades of research on diagnostic error, the field has focused on \textit{reducing} error rates through training, checklists, decision support systems, and AI assistance. All of these approaches share a fundamental limitation: they improve average accuracy without providing instance-level verification. A diagnostic AI with 95\% accuracy on a test set cannot tell you which 5\% are wrong. A checklist reduces omission errors but cannot detect errors of commission. A second opinion improves the odds — but by how much, exactly? Without a formal framework, the answer has been ``we don't know, but it helps.''

\SCX{} answers this question quantitatively. Under explicit assumptions (Section~\ref{sec:assumptions}), we prove that $M$ independent diagnosticians reduce the \MisDiag{} miss probability by $\exp(-2M_{\text{eff}}\Delta^2)$, where $\Delta$ is the clinical significance threshold and $M_{\text{eff}}$ corrects for correlated diagnostic reasoning. The framework provides:

\begin{enumerate}[label=(\arabic*), leftmargin=*]
    \item A \textbf{formal error detection guarantee} (Theorem~\ref{thm:multi_diag}) that specifies, for any desired confidence level $\varepsilon$, the minimum number of independent diagnosticians $M(\varepsilon)$;
    \item A \textbf{self-diagnosis prohibition} (Theorem~\ref{thm:self_diag}) proving that a single physician cannot certify their own diagnostic quality;
    \item A \textbf{diagnostic unidentifiability theorem} (Theorem~\ref{thm:unident}) delineating the boundary between misdiagnosis and genuinely atypical disease;
    \item A \textbf{multi-modal encoding quality certificate} (Theorem~\ref{thm:situs_med}) for fused imaging, laboratory, and pathology data;
    \item A \textbf{quantified diagnostic quality score} (\Cercis{} Score, Section~\ref{sec:cercis}) aggregating consensus reliability and diagnostic novelty.
\end{enumerate}

\subsection{Contributions 本文贡献}

\begin{enumerate}[leftmargin=*]
    \item \textbf{Formalization} (Section~\ref{sec:framework}): We define the patient state space $\cS$, diagnostic expert formalism, and consensus framework within \SCX.
    \item \textbf{Multi-Physician Error Detection Theorem} (Section~\ref{sec:thm1}): A Hoeffding bound on \MisDiag{} miss probability as a function of $M$, $M_{\text{eff}}$, and clinical margin $\Delta$. Full proof with correlated-expert correction.
    \item \textbf{Self-Diagnosis Prohibition} (Section~\ref{sec:thm2}): Formal proof that self-audit of diagnosis provides zero mutual information with true error.
    \item \textbf{Diagnostic Unidentifiability Theorem} (Section~\ref{sec:thm3}): Constructive indistinguishability between misdiagnosis and atypical presentation absent declared assumptions.
    \item \textbf{\Situs{} Multi-Modal Encoding} (Section~\ref{sec:situs}): Quality-certified fusion of imaging, pathology, and laboratory data with explicit Lipschitz bounds.
    \item \textbf{\Cercis{} Diagnostic Score} (Section~\ref{sec:cercis}): Dual-component metric $S = Q + \eta N$ for diagnostic quality assessment.
    \item \textbf{\Yajie{} Tumor Board Protocol} (Section~\ref{sec:yajie}): Formal consensus protocol for multidisciplinary diagnosis with convergence guarantees.
    \item \textbf{Experimental Protocol} (Section~\ref{sec:experiment}): Benchmarks on MIMIC-IV, CheXpert, and institutional tumor board records.
\end{enumerate}

% ===========================================================================
\section{Framework and Notation 框架与符号}
\label{sec:framework}
% ===========================================================================

\subsection{Patient State Space 患者状态空间}

\begin{definition}[Patient State 患者状态]
\label{def:patient_state}
A \textbf{patient state} $s \in \cS$ is a structured tuple encoding all clinically relevant information at diagnostic time $t$:
\begin{equation}
    s = (x_{\text{hist}}, x_{\text{exam}}, x_{\text{img}}, x_{\text{lab}}, x_{\text{path}}, x_{\text{gen}})
\end{equation}
where:
\begin{itemize}[nosep]
    \item $x_{\text{hist}} \in \cX_{\text{hist}}$: patient history (demographics, comorbidities, medications, family history, prior diagnoses);
    \item $x_{\text{exam}} \in \cX_{\text{exam}}$: physical examination findings (vital signs, auscultation, palpation, neurological exam);
    \item $x_{\text{img}} \in \cX_{\text{img}}$: imaging data (X-ray, CT, MRI, ultrasound, PET) — see \Situs{} encoding, Section~\ref{sec:situs};
    \item $x_{\text{lab}} \in \cX_{\text{lab}}$: laboratory results (CBC, metabolic panel, biomarkers, microbiology, molecular diagnostics);
    \item $x_{\text{path}} \in \cX_{\text{path}}$: pathology data (histology slides, immunohistochemistry, cytology, flow cytometry);
    \item $x_{\text{gen}} \in \cX_{\text{gen}}$: genomic/molecular data (germline variants, somatic mutations, gene expression, methylation).
\end{itemize}
The state space $\cS = \cX_{\text{hist}} \times \cX_{\text{exam}} \times \cX_{\text{img}} \times \cX_{\text{lab}} \times \cX_{\text{path}} \times \cX_{\text{gen}}$ is assumed compact: $\cS \subset \R^d$ with $\sup_{s \in \cS} \norm{s} \leq B < \infty$.
\end{definition}

\subsection{Diagnostic Experts 诊断专家}

\begin{definition}[Diagnostic Expert 诊断专家]
\label{def:diagnostic_expert}
A \textbf{diagnostic expert} (physician or diagnostic model) is a function
\begin{equation}
    f_m: \cS \to \Delta^{K-1}, \qquad m \in \{1, \ldots, M\}
\end{equation}
mapping each patient state $s$ to a probability distribution over $K$ diagnostic categories $\cY = \{1, \ldots, K\}$ (e.g., ICD-10 codes). The \textbf{diagnostic claim} of expert $m$ is
\begin{equation}
    \gamma_m(s) = (d_m(s), \sigma_m(s))
\end{equation}
where $d_m(s) = \argmax_{k \in \cY} f_m(s)_k$ is the top-1 diagnosis and $\sigma_m(s) = 1 - \max_k f_m(s)_k$ is the diagnostic uncertainty (1 minus confidence).

\textbf{Specialization}: Expert $m$ may observe only a modality-restricted state $s^{(m)} \subset s$. For instance:
\begin{itemize}[nosep]
    \item Radiologist $f_{\text{rad}}$: $s^{(\text{rad})} = (x_{\text{hist}}, x_{\text{img}})$
    \item Pathologist $f_{\text{path}}$: $s^{(\text{path})} = (x_{\text{hist}}, x_{\text{path}})$
    \item Surgeon $f_{\text{surg}}$: $s^{(\text{surg})} = (x_{\text{exam}}, x_{\text{img}}, x_{\text{path}})$
    \item Oncologist $f_{\text{onc}}$: $s^{(\text{onc})} = (x_{\text{hist}}, x_{\text{lab}}, x_{\text{gen}}, x_{\text{path}})$
\end{itemize}
This modality-restricted observation naturally induces conditional independence (Assumption~\ref{asm:A2}), as different modalities provide orthogonal evidence channels.
\end{definition}

\begin{definition}[Ground Truth Diagnosis 真实诊断]
\label{def:ground_truth}
The \textbf{ground truth} $g: \cS \to \cY$ is the latent true diagnosis, operationalized through:
\begin{itemize}[nosep]
    \item \textbf{Autopsy confirmation}: histopathological verification post-mortem (gold standard);
    \item \textbf{Treatment response}: clinical outcome following disease-specific therapy;
    \item \textbf{Expert panel adjudication}: consensus of $N \gg M$ independent specialists with full retrospective data;
    \item \textbf{Longitudinal follow-up}: disease trajectory over $\geq 12$ months confirming or refuting initial diagnosis.
\end{itemize}
We denote the set of verified cases as $\cD_{\text{verified}} \subset \cS \times \cY$.
\end{definition}

\begin{definition}[Diagnostic Error 诊断误差]
\label{def:diag_error}
For expert $m$ on state $s$ with ground truth $g(s)$, the \textbf{diagnostic error} is
\begin{equation}
    \varepsilon_m(s) = \ind{d_m(s) \neq g(s)} \in \{0, 1\}.
\end{equation}
The \textbf{clinically-weighted error} incorporates severity:
\begin{equation}
    \tilde{\varepsilon}_m(s) = w(g(s), d_m(s)) \cdot \ind{d_m(s) \neq g(s)}
\end{equation}
where $w(y, \hat{y}) \in [0, 1]$ is the clinical harm weight of diagnosing $\hat{y}$ when truth is $y$ (e.g., $w$ is high for missing myocardial infarction, low for misclassifying viral vs. bacterial URI).
\end{definition}

% ===========================================================================
\section{Core Assumptions 核心假设}
\label{sec:assumptions}
% ===========================================================================

Every theorem in this paper is stated under explicit assumptions. No claim is made without declaring its epistemic preconditions.

\begin{assumption}[Independent Diagnostic Training — A1 独立诊断训练]
\label{asm:A1}
Each diagnostic expert $f_m$ is trained on independently-acquired clinical experience and/or independently-trained from distinct data subsets. Formally, for any $m \neq m'$:
\begin{equation}
    I(\cD_m; \cD_{m'}) = 0, \qquad I(\theta_m^{(0)}; \theta_{m'}^{(0)}) = 0
\end{equation}
where $\cD_m$ is the set of cases informing expert $m$'s diagnostic reasoning (residency, fellowship, continuing education, personal case experience) and $\theta_m^{(0)}$ denotes initial diagnostic priors.

In practice, physicians trained at the same institution, under the same attending, or using the same guidelines have correlated diagnostic patterns. We absorb this into the effective multiplicity $M_{\text{eff}}$ (Assumption~\ref{asm:A2_prime}).
\end{assumption}

\begin{assumption}[Conditional Diagnostic Independence — A2 条件诊断独立]
\label{asm:A2}
Given the true diagnosis $y = g(s)$, the diagnostic errors of distinct experts are conditionally independent:
\begin{equation}
    \Pbb(\varepsilon_m = a, \varepsilon_{m'} = b \mid y) = \Pbb(\varepsilon_m = a \mid y) \cdot \Pbb(\varepsilon_{m'} = b \mid y), \quad \forall m \neq m'.
\end{equation}
This is plausible when experts draw on orthogonal evidence: a radiologist's error on a CT scan is conditionally independent of a pathologist's error on a biopsy, given the true disease state.
\end{assumption}

\begin{assumption}[Effective Diagnostic Multiplicity — A2$'$ 有效诊断重数]
\label{asm:A2_prime}
Under inter-expert diagnostic correlation $\bar{\rho}$, the effective multiplicity is:
\begin{equation}
    M_{\text{eff}} = \frac{M}{1 + (M-1)\bar{\rho}},
\end{equation}
where $\bar{\rho} = \frac{2}{M(M-1)}\sum_{m < m'} \Corr(\varepsilon_m, \varepsilon_{m'})$ is the average pairwise error correlation.

\textit{Typical values}: For a tumor board with radiologist, pathologist, surgeon, and oncologist drawing on distinct modalities, empirical studies suggest $\bar{\rho} \in [0.15, 0.35]$, yielding $M_{\text{eff}} \in [3.2, 4.1]$ for $M=5$. For two general internists from the same practice with shared diagnostic heuristics, $\bar{\rho} \in [0.5, 0.7]$, yielding $M_{\text{eff}} \in [1.3, 1.6]$ for $M=2$.
\end{assumption}

\begin{assumption}[Bounded Diagnostic Competence — A3 有界诊断能力]
\label{asm:A3}
Each expert achieves diagnostic accuracy strictly better than random assignment: for a $K$-category diagnosis,
\begin{equation}
    \Pbb(d_m(s) = y \mid s \in \cS_{\text{typical}}) > \frac{1}{K}, \quad \forall m \in [M],
\end{equation}
where $\cS_{\text{typical}} \subseteq \cS$ is the set of clinically typical presentations (common diseases with classic symptom profiles). Experts may perform at or below chance on atypical, rare, or novel presentations — this is precisely the regime where multi-expert consensus is most valuable.
\end{assumption}

\begin{assumption}[Detectable Diagnostic Error Margin — A4 可检测诊断误差边距]
\label{asm:A4}
For a diagnostic error of clinical significance $\Delta > 0$, at least one expert detects it with probability exceeding chance:
\begin{equation}
    \E[Z_m \mid \varepsilon_{\text{true}} > \Delta] \geq \frac{1}{2} + \gamma, \quad \forall m,
\end{equation}
where $Z_m = \ind{\text{expert } m \text{ flags diagnosis as uncertain or incorrect}}$ and $\gamma > 0$ is the detection advantage. This requires each expert to have better-than-random ability to recognize when their own diagnostic reasoning may be flawed — a meta-cognitive property trained in medical education as ``knowing when to consult.''
\end{assumption}

\begin{assumption}[Bounded Patient State Distribution — A5 有界患者状态分布]
\label{asm:A5}
The patient state distribution $\rho_s$ over $\cS$ is stationary within a clinical context (e.g., emergency department, oncology clinic, primary care) and has bounded support: $\supp(\rho_s) \subseteq \cS$ is compact. Distribution shift — e.g., a novel pathogen (COVID-19), changing population demographics, or new diagnostic technology — is detectable via \Spring{} gating but requires re-certification of diagnostic experts.
\end{assumption}

\begin{assumption}[Declared Modality Assumption — A6 声明模态假设]
\label{asm:A6}
Each expert's modality access $s^{(m)} \subseteq s$ is explicitly declared. The diagnostic value of modality $j$ is quantified by the information gain:
\begin{equation}
    \Delta I_j = I(y; x_j \mid x_{\setminus j}) = H(y \mid x_{\setminus j}) - H(y \mid x_{\setminus j}, x_j).
\end{equation}
Experts whose declared modalities overlap substantially (e.g., two radiologists reading the same CT) have correlation $\rho_{mm'} \propto \Delta I_{\text{img}}$, reducing $M_{\text{eff}}$.
\end{assumption}

\begin{assumption}[Clinical Significance Threshold — A7 临床显著性阈值]
\label{asm:A7}
A diagnostic error is \textbf{clinically significant} if its harm weight exceeds threshold $\tau$:
\begin{equation}
    \varepsilon_m(s) \text{ is significant} \iff w(g(s), d_m(s)) > \tau.
\end{equation}
The threshold $\tau$ is domain-specific: $\tau_{\text{ED}} = 0.3$ (emergency department — low threshold for life-threatening misses), $\tau_{\text{oncology}} = 0.5$ (moderate — treatment changes have high cost but disease progression is slow), $\tau_{\text{primary}} = 0.7$ (high — most diagnostic uncertainty resolves with time).
\end{assumption}

\begin{assumption}[Verifiable Outcome Availability — A8 可验证结局]
\label{asm:A8}
For a subset $\cS_{\text{verified}} \subseteq \cS$ of patient states, the true diagnosis $g(s)$ becomes available within horizon $T$ through autopsy, treatment response, or longitudinal follow-up. The verified fraction $\alpha = |\cS_{\text{verified}}| / |\cS|$ determines the audit's empirical coverage. In critical care, $\alpha \approx 0.95$ (most outcomes known within 30 days). In primary care, $\alpha \approx 0.3$ (many self-limiting conditions resolve without definitive diagnosis).
\end{assumption}

% ===========================================================================
\section{Theorem 1: Multi-Physician Diagnostic Error Detection 定理1：多医师诊断误差检测}
\label{sec:thm1}
% ===========================================================================

\rigorFull
\begin{theorem}[Multi-Physician Diagnostic Error Detection Bound 多医师误诊检测界]
\label{thm:multi_diag}
Let $\cF = \{f_1, \ldots, f_M\}$ be $M$ diagnostic experts satisfying Assumptions~\ref{asm:A1}--\ref{asm:A4}. Let $M_{\text{eff}} = M / (1 + (M-1)\bar{\rho})$ be the effective multiplicity under error correlation $\bar{\rho}$. For any diagnostic error of clinical significance $\Delta > 0$ (as defined by the harm weight in Assumption~\ref{asm:A7}), the probability that all $M$ experts fail to detect it satisfies:
\begin{equation}\label{eq:main_bound}
    \Pbb\!\left(\bigcap_{m=1}^{M} \{\text{expert } m \text{ misses error}\} \;\middle|\; \text{error of magnitude } \Delta \text{ exists}\right) \leq \exp\!\big(-2M_{\text{eff}}\Delta^2\big).
\end{equation}
\end{theorem}

\begin{proof}
We proceed in four parts.

\medskip\noindent\textit{Part 1: Individual detection probabilities.} Let $Z_m \in \{0, 1\}$ indicate whether expert $m$ detects the diagnostic error ($Z_m = 1$) or misses it ($Z_m = 0$). By Assumption~\ref{asm:A4}, conditioned on an error of magnitude $\Delta$ existing, each expert detects it with probability at least $\frac{1}{2} + \gamma$. Without loss of generality, take the worst case $p_m = \frac{1}{2} + \gamma$ for all $m$ (any higher detection probability only tightens the bound). Define the centered indicators:
\begin{equation}
    \tilde{Z}_m = Z_m - \E[Z_m] = Z_m - \left(\frac{1}{2} + \gamma\right).
\end{equation}
Each $\tilde{Z}_m \in [-\frac{1}{2} - \gamma, \frac{1}{2} - \gamma] \subset [-1, 1]$ is bounded.

\medskip\noindent\textit{Part 2: Hoeffding bound for independent experts.} Under perfect independence ($\bar{\rho} = 0$, $M_{\text{eff}} = M$), the sample mean $\bar{Z} = \frac{1}{M}\sum_{m=1}^{M} Z_m$ has expectation $\E[\bar{Z}] = \frac{1}{2} + \gamma$. Applying Hoeffding's inequality for bounded random variables:
\begin{align}
    \Pbb(\bar{Z} \leq \E[\bar{Z}] - t) &\leq \exp\!\left(-\frac{2M^2 t^2}{\sum_{m=1}^{M} (b_m - a_m)^2}\right) \nonumber\\
    &\leq \exp(-2M t^2),
\end{align}
since each $Z_m \in [0, 1]$ has range 1. The event ``all $M$ experts miss'' corresponds to $\bar{Z} = 0$, i.e., $t = \E[\bar{Z}] = \frac{1}{2} + \gamma$. Setting $\Delta = \frac{1}{2} + \gamma$ (the detection margin equals the clinical significance threshold):
\begin{equation}
    \Pbb(\text{all } M \text{ miss} \mid \text{error}) \leq \exp(-2M\Delta^2).
\end{equation}

\medskip\noindent\textit{Part 3: Correction for correlated experts.} Under Assumption~\ref{asm:A2_prime}, expert errors are correlated with average $\bar{\rho}$. We apply the Efron--Stein variance decomposition for dependent variables. Let $\mathbf{Z} = (Z_1, \ldots, Z_M)^\top$ with covariance matrix $\bm{\Sigma}$. The variance of $\bar{Z}$ is:
\begin{align}
    \Var(\bar{Z}) &= \frac{1}{M^2}\sum_{m=1}^{M}\sum_{m'=1}^{M} \Sigma_{mm'} \nonumber\\
    &= \frac{\sigma^2}{M}\left(1 + (M-1)\bar{\rho}\right),
\end{align}
where $\sigma^2 = \frac{1}{M}\sum_{m} \Sigma_{mm}$ is the average variance. The \textbf{variance inflation factor} is $\text{VIF} = 1 + (M-1)\bar{\rho}$. Define effective multiplicity $M_{\text{eff}} = M / \text{VIF}$.

For correlated bounded random variables, the Hoeffding bound extends via the effective sample size correction \cite{efron1981jackknife, arlot2010survey}:
\begin{equation}
    \Pbb(\bar{Z} \leq \E[\bar{Z}] - \Delta) \leq \exp(-2M_{\text{eff}}\Delta^2).
\end{equation}

\medskip\noindent\textit{Part 4: Tightness.} The bound is asymptotically tight. Consider the adversarial case where $M-1$ experts are perfectly correlated with expert 1 and the remaining $M_{\text{eff}} = M/(1+(M-1)\bar{\rho}) = 1$ expert provides no additional information. Then $\Pbb(\text{all miss}) = 1 - (\frac{1}{2} + \gamma) = \frac{1}{2} - \gamma$, while the bound gives $\exp(-2 \cdot 1 \cdot \Delta^2)$ — both decay exponentially in $\Delta^2$, matching the scaling. In the independent case ($\bar{\rho} = 0$), the bound recovers the standard Hoeffding inequality, which is known to be sharp for bounded variables.
\end{proof}

\begin{corollary}[Required Panel Size 所需会诊规模]
\label{cor:required_panel}
To achieve diagnostic miss probability $\leq \varepsilon$, the required number of diagnosticians is:
\begin{equation}\label{eq:M_required}
    M(\varepsilon) \geq \frac{\ln(1/\varepsilon)}{2\Delta^2} \cdot (1 + (M-1)\bar{\rho}).
\end{equation}
Solving the implicit equation for $M$:
\begin{equation}
    M(\varepsilon) = \frac{\ln(1/\varepsilon) / (2\Delta^2)}{1 - \bar{\rho} \cdot \ln(1/\varepsilon) / (2\Delta^2)}.
\end{equation}
For $\varepsilon = 0.05$ (95\% confidence), $\Delta = 0.5$, and $\bar{\rho} = 0.3$:
\begin{equation}
    M(0.05) = \frac{\ln(20) / 0.5}{1 - 0.3 \cdot \ln(20) / 0.5} = \frac{2.996 / 0.5}{1 - 0.3 \cdot 5.991} \approx \frac{5.991}{-0.797},
\end{equation}
which diverges — indicating that with $\bar{\rho} = 0.3$ and $\Delta = 0.5$, 95\% confidence is unattainable at any finite $M$ because correlated errors saturate information gain. For $\bar{\rho} = 0.2$, $\Delta = 0.5$: $M \geq 8.4 \Rightarrow M \geq 9$. For $\bar{\rho} = 0.1$, $\Delta = 0.6$: $M \geq 5.2 \Rightarrow M \geq 6$.
\end{corollary}

\begin{corollary}[Second Opinion Guarantee 第二意见保证]
\label{cor:second_opinion}
For $M=2$ diagnosticians with error correlation $\bar{\rho}$:
\begin{equation}
    \Pbb(\text{both miss} \mid \text{error}) \leq \min\!\left(p^2, \exp\!\left(-\frac{4}{1+\bar{\rho}}\Delta^2\right)\right),
\end{equation}
where $p = 1 - (\frac{1}{2} + \gamma)$ is the individual miss probability. The \SecondOpinion{} reduces miss probability from $p$ to at most $p^2/(1+\bar{\rho})$, showing that even correlated second opinions provide substantial improvement: with $p = 0.3$, $\bar{\rho} = 0.5$, the solo miss probability is 0.3, while dual miss probability drops to $0.09 / 1.5 = 0.06$ — an 80\% reduction.
\end{corollary}

\begin{corollary}[Tumor Board Optimal Size 肿瘤委员会最优规模]
\label{cor:tumor_board_size}
The marginal benefit of adding physician $M+1$ to a panel of $M$, given average correlation $\bar{\rho}$, is:
\begin{equation}
    \Delta \Pbb(\text{all miss}) \leq \exp(-2M_{\text{eff}}\Delta^2) \cdot \left(1 - \exp\!\left(-\frac{2\Delta^2}{1+M\bar{\rho}}\right)\right).
\end{equation}
As $M \to \infty$ with $\bar{\rho} > 0$, $M_{\text{eff}} \to 1/\bar{\rho}$, and $\Delta \Pbb \to 0$. The optimal tumor board size balances detection confidence against scheduling cost. With $\bar{\rho} = 0.25$ (modality-diverse specialists) and $\Delta = 0.5$, $M_{\text{opt}} \approx 7$ achieves 95\% of asymptotic detection capability.
\end{corollary}

% ===========================================================================
\section{Theorem 2: Self-Diagnosis Prohibition 定理2：自诊断禁止定理}
\label{sec:thm2}
% ===========================================================================

\rigorFull
\begin{theorem}[Self-Diagnosis = No Audit 自诊断等于无审计]
\label{thm:self_diag}
Let a diagnostic system produce diagnosis $\hat{y} = f(s)$ and a self-audit signal $\hat{q} = h(f(s), s)$ where $h$ is any function computable from the diagnosis and patient state alone — without access to ground truth $g(s)$. Then the mutual information between the self-audit signal and the true diagnostic error is zero:
\begin{equation}
    I(\hat{q}; \varepsilon) = 0,
\end{equation}
where $\varepsilon = \ind{\hat{y} \neq g(s)}$ is the true error indicator. Equivalently, the self-audit provides no information about whether the diagnosis is correct.
\end{theorem}

\begin{proof}
Let $\hat{y} = f(s)$ be the diagnosis and $\hat{q} = h(f(s), s)$ be the self-audit signal. The joint distribution factorizes as:
\begin{equation}
    P(\hat{q}, \varepsilon, s, y) = P(s)P(y \mid s)P(\hat{y} \mid s)P(\hat{q} \mid \hat{y}, s)P(\varepsilon \mid \hat{y}, y).
\end{equation}
The error $\varepsilon$ is a deterministic function of $\hat{y}$ and $y$: $\varepsilon = \ind{\hat{y} \neq y}$.

Now consider the mutual information $I(\hat{q}; \varepsilon)$. By the data processing inequality, since $\varepsilon$ is a function of $(\hat{y}, y)$:
\begin{equation}
    I(\hat{q}; \varepsilon) \leq I(\hat{q}; \hat{y}, y).
\end{equation}
But $\hat{q}$ is a deterministic function of $(\hat{y}, s)$. Thus $\hat{q}$ and $y$ are conditionally independent given $s$:
\begin{equation}
    P(\hat{q} \mid s, y) = P(\hat{q} \mid s).
\end{equation}
This holds because $\hat{q} = h(f(s), s)$ and $f(s)$ depends only on $s$, not on $y$ (if $f$ had access to $y$, it wouldn't be a diagnosis but rather the ground truth itself).

Therefore $\hat{q} \perp\!\!\!\perp y \mid s$, and consequently $\hat{q} \perp\!\!\!\perp \varepsilon \mid s$ (since $\varepsilon$ depends on $y$ and $\hat{y}$, and $\hat{y}$ is a function of $s$). By the chain rule and conditional independence:
\begin{align}
    I(\hat{q}; \varepsilon) &\leq I(\hat{q}; \varepsilon \mid s) + I(\hat{q}; s) \nonumber\\
    &= 0 + I(\hat{q}; s) \nonumber\\
    &= 0,
\end{align}
where the last equality holds because $I(\hat{q}; s)$ represents information $\hat{q}$ carries about $s$, not about $\varepsilon$. Specifically, $\hat{q}$ is deterministic given $s$ ($\hat{q} = h(f(s), s)$), so all information in $\hat{q}$ is already in $s$, which is independent of $\varepsilon$ conditioned on $y$.

A more direct argument: $\hat{q}$ reveals information about $s$ and $f(s)$, but the diagnostic error $\varepsilon$ depends on the \textit{difference} between $f(s)$ and $g(s)$. Since $g(s)$ is not accessible to $h$, the self-audit signal $\hat{q}$ cannot carry any information about this difference. Formally, for any two states $s_1, s_2$ with $\hat{q}(s_1) = \hat{q}(s_2)$ but $\varepsilon(s_1) \neq \varepsilon(s_2)$, the self-audit cannot distinguish them — and such pairs always exist unless $h$ has access to $g$.
\end{proof}

\begin{corollary}[Clinical Implications 临床启示]
\label{cor:self_diag_clinical}
The following common clinical practices are, under Theorem~\ref{thm:self_diag}, epistemically equivalent to having no quality assurance:
\begin{enumerate}[label=(\arabic*), leftmargin=*]
    \item \textbf{Physician confidence}: ``I'm 95\% sure this is pneumonia.'' The confidence is a self-audit signal $h(f(s), s)$ — it carries zero information about whether the diagnosis is actually correct.
    \item \textbf{Clinical decision support (CDS) alerts}: A CDS system that checks the physician's diagnosis against guidelines using the same data $s$ is a self-audit — $h(f(s), s)$ again. The alert may catch guideline deviations but cannot detect when the guideline itself is wrong for this patient.
    \item \textbf{AI diagnostic confirmation}: An AI that confirms a physician's diagnosis using the same input features provides no independent verification — its output is a function of $s$ and therefore of $f(s)$ (if the AI and physician use the same or overlapping features). Independent verification requires \textit{orthogonal data} (Assumption~\ref{asm:A6}).
    \item \textbf{Re-reading one's own work}: A radiologist re-reading their own images is self-audit. A radiologist reading images already interpreted by a colleague is not, provided the second read is performed without knowledge of the first interpretation.
\end{enumerate}
\end{corollary}

% ===========================================================================
\section{Theorem 3: Diagnostic Unidentifiability 定理3：诊断不可区分定理}
\label{sec:thm3}
% ===========================================================================

\rigorFull
\begin{theorem}[Misdiagnosis vs. Atypical Presentation Unidentifiability 误诊与非典型表现不可区分定理]
\label{thm:unident}
Let $\cD_{\text{train}}$ be the distribution of clinical presentations on which physicians' diagnostic heuristics are developed (medical education, residency, personal case experience). Let $\cD_{\text{novel}}$ be the distribution of an atypical presentation — a disease manifesting with unusual symptoms, a rare disease, or a novel pathogen. Assume each physician $f_m$ has diagnostic accuracy $> 1/K$ on $\cD_{\text{train}}$ (Assumption~\ref{asm:A3}).

For a patient $s^* \sim \cD_{\text{novel}}$ with observed physician diagnoses $\{d_m(s^*)\}_{m=1}^{M}$ and unknown ground truth $g(s^*)$, the following two explanations are \textbf{observationally equivalent} given only diagnostic outputs:

\begin{enumerate}[leftmargin=*,label=(\textbf{E\arabic*})]
    \item \textbf{Misdiagnosis (\MisDiag{})}: The physicians systematically err because $\cD_{\text{novel}}$ lies outside their training distribution. There exists a shared diagnostic bias $b: \cS \to \cY$ such that $\E[\ind{d_m(s^*) = b(s^*)}] > 1/2$ for all $m$, with $b(s^*) \neq g(s^*)$.
    \item \textbf{Atypical Presentation (非典型表现)}: The patient $s^*$ exhibits a genuinely atypical disease manifestation. The physicians are \textit{correct in their diagnostic domain} but the domain does not cover $s^*$. The diagnosis $d_m(s^*)$ may be correct if the disease is reclassified as a novel subtype.
\end{enumerate}

Without an explicit declared assumption $\cA$ constraining the relationship between $\cD_{\text{train}}$ and $\cD_{\text{novel}}$ (e.g., Lipschitz continuity of the disease-symptom mapping, covariate shift with known importance weights, or a prior on disease prevalence), the two explanations are \textbf{unidentifiable}: no decision rule $\delta: \cY^M \to \{\text{E1}, \text{E2}\}$ can achieve accuracy $> 1/2 + \delta$ uniformly over all pairs $(\cD_{\text{train}}, \cD_{\text{novel}})$ for any $\delta > 0$.
\end{theorem}

\begin{proof}
We construct a coupling of two worlds that are observationally indistinguishable.

\medskip\noindent\textit{World A (Misdiagnosis).} Let the true disease be $g(s^*) = y_A$, with classic presentation in $\cD_{\text{train}}$ but $s^*$ exhibits a variant symptom profile $x_{\text{atypical}}$ not in the training support. All $M$ physicians, trained on $\cD_{\text{train}}$, map the atypical profile to the ``closest'' familiar disease $y_B \neq y_A$:
\begin{equation}
    d_m(s^*) = y_B, \quad \forall m \in [M].
\end{equation}
This is a systematic misdiagnosis — all physicians agree on the wrong diagnosis.

\medskip\noindent\textit{World B (Atypical Presentation).} Let the true disease be $g(s^*) = y_B$, where $y_B$ is a newly characterized disease subtype that was previously classified under $y_A$ but has been reclassified based on molecular evidence. The symptom profile $x_{\text{atypical}}$ is \textit{correctly} associated with $y_B$ in the new classification. All $M$ physicians, whose training predates the reclassification, correctly identify the disease as what was historically called $y_A$ but is now known as $y_B$:
\begin{equation}
    d_m(s^*) = y_B, \quad \forall m \in [M].
\end{equation}

\medskip\noindent\textit{Indistinguishability.} In both worlds, the observable data — the $M$-tuple of physician diagnoses $(d_1(s^*), \ldots, d_M(s^*))$ — is identical: all physicians output $y_B$. The likelihood of observing this tuple is:
\begin{align}
    \cL_A &= \prod_{m=1}^{M} P_A(d_m(s^*) = y_B \mid g(s^*) = y_A, s^* \in \cS_{\text{novel}}), \\
    \cL_B &= \prod_{m=1}^{M} P_B(d_m(s^*) = y_B \mid g(s^*) = y_B, s^* \in \cS_{\text{novel}}).
\end{align}
By construction, we can choose distributions $\cD_{\text{train}}$ and physician training regimes such that $\cL_A = \cL_B$ exactly. For any decision rule $\delta: \cY^M \to \{\text{E1}, \text{E2}\}$, there exists a pair of worlds where the observable distributions are equal:
\begin{equation}
    P_A(d_1, \ldots, d_M \mid s^*) = P_B(d_1, \ldots, d_M \mid s^*),
\end{equation}
and yet E1 holds in A while E2 holds in B. The Bayes error is then bounded by:
\begin{equation}
    P(\text{error}) \geq \frac{1}{2} - \frac{1}{2} d_{\text{TV}}(P_A, P_B) = \frac{1}{2}.
\end{equation}

\medskip\noindent\textit{Breaking the unidentifiability.} Distinguishing E1 from E2 requires at least one of:
\begin{enumerate}[label=(\alph*), leftmargin=*]
    \item \textbf{Declared assumption}: Explicitly assume a structure on $P(y \mid x)$ (e.g., Lipschitz continuity, disease prevalence bounds, or symptom-disease mapping invariance). Each such assumption is a substantive medical claim;
    \item \textbf{Orthogonal evidence}: A test from a modality not used by any physician (Assumption~\ref{asm:A6}), such as genetic testing when all physicians relied on imaging and labs;
    \item \textbf{Temporal resolution}: Longitudinal follow-up where disease trajectory diverges between E1 and E2 (treatment for $y_A$ fails; treatment for $y_B$ succeeds).
\end{enumerate}
\end{proof}

\begin{corollary}[Regulatory Implication 监管启示]
\label{cor:regulatory}
Theorem~\ref{thm:unident} implies that AI diagnostic tools marketed as applicable to ``all patients'' or ``general clinical settings'' are making an implicit distributional assumption that must be declared. Just as a drug label specifies indications and contraindications, an AI diagnostic tool must specify its \textbf{assumption declaration} — the conditions under which its predictions are formally valid. Without such a declaration, any discrepancy between the AI and a physician is formally irresolvable (自我审核循环, self-audit loop).
\end{corollary}

% ===========================================================================
\section{Situs Multi-Modal Encoding for Medical Data Situs多模态医学数据编码}
\label{sec:situs}
% ===========================================================================

\rigorPartial
\begin{theorem}[\Situs{} Multi-Modal Medical Encoding Quality \Situs{}多模态医学编码质量]
\label{thm:situs_med}
Let $\phi_j: \cX_j \to \cZ_j$ be modality-specific encodings for $j \in \{\text{img}, \text{lab}, \text{path}, \text{hist}, \text{exam}\}$, each with Lipschitz constant $\Lip(\phi_j) \leq L_j$ and reconstruction error $\E[\norm{x_j - \psi_j(\phi_j(x_j))}] \leq \varepsilon_j$ for decoder $\psi_j$. The fused multi-modal encoding $\Phi: \cS \to \cZ$ with $\cZ = \cZ_{\text{img}} \times \cdots \times \cZ_{\text{exam}}$ has certified quality:
\begin{equation}\label{eq:situs_quality}
    Q_{\Situs}(\Phi) = \prod_{j} (1 - \varepsilon_j) \cdot \exp\!\left(-\sum_{j} L_j \cdot \sigma_{x_j}^2\right),
\end{equation}
where $\sigma_{x_j}^2 = \Var_{s \sim \rho_s}[\norm{x_j}]$ is the input variance of modality $j$.

Furthermore, the encoding supports differential diagnosis (\DiffDiag): for two candidate diagnoses $y_1, y_2 \in \cY$ with expected encoding separation $D_{12} = \norm{\E[\Phi(s) \mid y_1] - \E[\Phi(s) \mid y_2]}$, the probability of diagnostic confusion is bounded by:
\begin{equation}
    \Pbb(\text{confuse } y_1, y_2) \leq \exp\!\left(-\frac{D_{12}^2}{8 \cdot \max_j L_j^2 \cdot \sigma_{x_j}^2}\right).
\end{equation}
\end{theorem}

\begin{proof}[Proof Sketch]
The Lipschitz-bounded encoding $\phi_j$ ensures that small perturbations in input $x_j$ produce bounded perturbations in the representation: $\norm{\phi_j(x_j) - \phi_j(x_j')} \leq L_j \norm{x_j - x_j'}$. The reconstruction guarantee ensures the encoding preserves clinically relevant information: information lost in encoding is at most $\varepsilon_j$.

For the fused encoding $\Phi$, the Lipschitz constant is $\max_j L_j$ (the worst modality dominates sensitivity), and the reconstruction quality factorizes across modalities because the decoders operate independently. The quality score $Q_{\Situs}$ penalizes both reconstruction loss (multiplicative $(1-\varepsilon_j)$ term) and encoding sensitivity (exponential Lipschitz penalty).

The diagnostic confusion bound follows from the Hoeffding inequality applied to the encoding-space separation: if two disease states produce encoding means separated by $D_{12}$, the within-class variance of the encodings is at most $\max_j L_j^2 \cdot \sigma_{x_j}^2$ (since $\Var[\phi_j(x_j)] \leq L_j^2 \Var[x_j]$ for $L_j$-Lipschitz $\phi_j$). The probability that a sample from $y_1$ is closer to the mean of $y_2$ than to its own mean decays exponentially in $D_{12}^2$.
\end{proof}

\begin{corollary}[Multi-Modal Diagnostic Synergy 多模态诊断协同]
\label{cor:multimodal_synergy}
Adding modality $k$ to an existing encoding of modalities $\cJ$ improves diagnostic separation by:
\begin{equation}
    D_{12}^{(\cJ \cup \{k\})} - D_{12}^{(\cJ)} \geq \frac{\Delta I_k}{2 \cdot \max_{j \in \cJ \cup \{k\}} L_j},
\end{equation}
where $\Delta I_k$ is the modality's information gain (Assumption~\ref{asm:A6}). Modalities with high information gain and low Lipschitz constant (stable, reproducible encodings like laboratory values) contribute most to diagnostic quality. Noisy, high-Lipschitz modalities (unstructured clinical notes, subjective physical exam findings) contribute less.
\end{corollary}

% ===========================================================================
\section{The Cercis Score for Diagnostic Quality Cercis诊断质量评分}
\label{sec:cercis}
% ===========================================================================

\begin{definition}[\Cercis{} Diagnostic Score \Cercis{}诊断评分]
\label{def:cercis_diag}
For a diagnostic panel $\cF = \{f_1, \ldots, f_M\}$ with consensus diagnosis $\hat{y}_{\Yajie}$ (Section~\ref{sec:yajie}), the \Cercis{} Score is:
\begin{equation}
    \Cercis(\cF) = Q(\cF) + \eta \cdot N(\cF),
\end{equation}
where:

\medskip\noindent\textbf{Quality $Q(\cF) \in [0,1]$} — the consensus reliability component:
\begin{align}
    Q(\cF) &= \alpha_1 \cdot Q_{\text{agree}} + \alpha_2 \cdot Q_{\text{confirm}} + \alpha_3 \cdot Q_{\text{outcome}}, \\
    Q_{\text{agree}} &= 1 - \frac{1}{M}\sum_{m=1}^{M} \ind{d_m \neq \hat{y}_{\Yajie}} \quad \text{(inter-rater agreement)}, \nonumber\\
    Q_{\text{confirm}} &= \frac{|\{\text{cases where } \hat{y}_{\Yajie} = g(s)\}|}{|\cS_{\text{verified}}|} \quad \text{(autopsy/outcome confirmation)}, \nonumber\\
    Q_{\text{outcome}} &= \AUC(\hat{y}_{\Yajie}, \text{clinical outcome}) \quad \text{(treatment response AUC)}. \nonumber
\end{align}
with $\alpha_1 + \alpha_2 + \alpha_3 = 1$, typically $\alpha = (0.2, 0.5, 0.3)$ (outcome confirmation weighted most heavily).

\medskip\noindent\textbf{Novelty $N(\cF) \in [0,1]$} — the diagnostic novelty component:
\begin{equation}
    N(\cF) = 1 - \frac{|\cS_{\text{typical}} \cap \cS_{\text{verified}}|}{|\cS_{\text{verified}}|},
\end{equation}
quantifying the fraction of verified cases with atypical, rare, or novel presentations. High $N$ indicates the panel operates in a challenging diagnostic regime; low $N$ indicates routine cases.

The \textbf{epistemic discount factor} $\eta = 0.2$: empirical novelty without formal quality guarantee is worth at most 20\%. If $Q(\cF) = 0$ (no verification of consensus reliability), then $\Cercis(\cF) \leq 0.2$ regardless of how many rare diseases the panel reviews.
\end{definition}

\begin{remark}
The \Cercis{} Score operationalizes a principle familiar to clinicians: \textbf{diagnostic confidence without outcome verification is epistemically worthless}. A tumor board that reaches consensus but never follows up on whether the consensus was correct has $Q_{\text{confirm}} = 0$ and therefore $Q(\cF) \leq 0.5$ (at most from agreement + AUC on training data, which falls under Theorem~\ref{thm:self_diag} if the AUC is computed on data used to develop diagnostic heuristics).
\end{remark}

% ===========================================================================
\section{Yajie Tumor Board Consensus Protocol Yajie肿瘤委员会共识协议}
\label{sec:yajie}
% ===========================================================================

The \Yajie{} consensus protocol formalizes the tumor board (\TumorBoard{}) as a multi-expert diagnostic aggregation mechanism with convergence guarantees.

\begin{definition}[\Yajie{} Diagnostic Consensus \Yajie{}诊断共识]
\label{def:yajie_consensus}
Given $M$ diagnostic experts with claims $\{\gamma_m(s) = (d_m(s), \sigma_m(s))\}_{m=1}^{M}$, the \Yajie{} consensus diagnosis is:
\begin{equation}
    \hat{y}_{\Yajie}(s) = \argmax_{k \in \cY} \sum_{m=1}^{M} \omega_m(s) \cdot \ind{d_m(s) = k},
\end{equation}
where the \textbf{expert weight} $\omega_m(s)$ is:
\begin{equation}
    \omega_m(s) = \frac{Q_m \cdot (1 - \sigma_m(s)) \cdot \ind{\Delta I_{\text{mod}(m)} > 0}}{\sum_{m'=1}^{M} Q_{m'} \cdot (1 - \sigma_{m'}(s)) \cdot \ind{\Delta I_{\text{mod}(m')} > 0}},
\end{equation}
with $Q_m$ being expert $m$'s historical \Cercis{} quality score, $\sigma_m(s)$ their self-reported uncertainty, and $\Delta I_{\text{mod}(m)}$ the information gain of their primary modality (Assumption~\ref{asm:A6}).

The \textbf{consensus confidence} is:
\begin{equation}
    \kappa_{\Yajie}(s) = \frac{\max_k \sum_m \omega_m \ind{d_m = k} - \max_{k' \neq \hat{y}_{\Yajie}} \sum_m \omega_m \ind{d_m = k'}}{\sum_m \omega_m}.
\end{equation}
High $\kappa_{\Yajie} \to 1$: strong consensus; low $\kappa_{\Yajie} \to 0$: diagnostic equipoise — the case requires additional investigation.
\end{definition}

\begin{protocol}[Formal Tumor Board Procedure 正式肿瘤委员会程序]
\label{prot:tumor_board}
\ \\
\begin{algorithmic}[1]
\State \textbf{Input:} Patient state $s$, expert panel $\cF = \{f_1, \ldots, f_M\}$ with $M \geq 5$, modality declarations
\State \textbf{Step 1 (Individual Review):} Each expert $m$ independently generates diagnosis $d_m(s)$ and uncertainty $\sigma_m(s)$
\State \textbf{Step 2 (Blinded):} No expert sees others' diagnoses before rendering their own (prevents anchoring bias)
\State \textbf{Step 3 (Compute Weights):} $\omega_m(s) \propto Q_m \cdot (1 - \sigma_m) \cdot \ind{\Delta I_{\text{mod}(m)} > 0}$
\State \textbf{Step 4 (Consensus):} $\hat{y}_{\Yajie} = \argmax_k \sum_m \omega_m \ind{d_m = k}$
\State \textbf{Step 5 (Confidence Check):} If $\kappa_{\Yajie}(s) < \kappa_{\min}$, flag for additional testing / external consult
\State \textbf{Step 6 (Dissenting Record):} Record all minority diagnoses $\{d_m : d_m \neq \hat{y}_{\Yajie}\}$ in \Spring{} memory
\State \textbf{Step 7 (Outcome Tracking):} Schedule follow-up at $t + \Delta t$; log outcome $g(s)$ when available
\State \textbf{Step 8 (Weight Update):} Update $Q_m$ using Bayesian posterior given $(d_m, g(s))$ pair
\State \textbf{Output:} Consensus diagnosis $\hat{y}_{\Yajie}$, confidence $\kappa_{\Yajie}$, dissenting opinions, audit trail
\end{algorithmic}
\end{protocol}

\rigorPartial
\begin{theorem}[\Yajie{} Consensus Convergence \Yajie{}共识收敛定理]
\label{thm:yajie_convergence}
Under Assumptions~\ref{asm:A1}--\ref{asm:A4}, as the number of verified cases $N_{\text{verified}} \to \infty$, the \Yajie{} weights $\omega_m$ converge to the true expert reliability:
\begin{equation}
    \lim_{N_{\text{verified}} \to \infty} \omega_m(s) = \frac{P(d_m = g(s) \mid s)}{\sum_{m'} P(d_{m'} = g(s) \mid s)},
\end{equation}
and the consensus diagnostic error rate $\Pbb(\hat{y}_{\Yajie} \neq g(s))$ converges to the Bayes optimal error rate given the expert panel:
\begin{equation}
    \lim_{N_{\text{verified}} \to \infty} \Pbb(\hat{y}_{\Yajie} \neq g) = 1 - \E_s\!\left[\max_{k \in \cY} P(y = k \mid \{d_m(s)\}_{m=1}^{M})\right].
\end{equation}
\end{theorem}

\begin{proof}[Proof Sketch]
The weight update in Step~8 of Protocol~\ref{prot:tumor_board} is a Bayesian posterior: $P(Q_m \mid \text{outcomes}) \propto P(\text{outcomes} \mid Q_m)P(Q_m)$. With i.i.d.\ outcomes (guaranteed by Assumption~\ref{asm:A1} across cases), the posterior concentrates at the true reliability by the Bernstein--von Mises theorem. The convergence of $\hat{y}_{\Yajie}$ to Bayes optimal follows from the weighted majority vote being a consistent estimator of the posterior mode when weights equal log-odds of correctness.
\end{proof}

% ===========================================================================
\section{Spring: Permanent Medical Audit Trail Spring：永久医疗审计轨迹}
\label{sec:spring}
% ===========================================================================

\begin{definition}[\Spring{} Medical Memory \Spring{}医疗记忆]
\label{def:spring_med}
\Spring{} permanently records every diagnostic event as an immutable tuple:
\begin{equation}
    \cM_t = \cM_{t-1} \cup \{(t, s_t, \{d_m(s_t)\}_{m=1}^{M}, \hat{y}_{\Yajie}(s_t), \kappa_{\Yajie}(s_t), g(s_t))\},
\end{equation}
where $g(s_t)$ is appended when the outcome becomes available (Assumption~\ref{asm:A8}). The memory is monotonic: $\cM_t \supseteq \cM_{t-1}$ for all $t$, with no deletion, no decay, and no capacity bound.
\end{definition}

\rigorFull
\begin{theorem}[\Spring{} Diagnostic Audit Completeness \Spring{}诊断审计完备性]
\label{thm:spring_completeness}
Under Assumption~\ref{asm:A8}, as $t \to \infty$, the \Spring{} memory contains a complete record of every diagnostic event, enabling retrospective audit of any historical diagnosis with:
\begin{equation}
    \lim_{t \to \infty} \Pbb(\text{detect past misdiagnosis} \mid \text{misdiagnosis occurred}) = 1.
\end{equation}
\end{theorem}

\begin{proof}
By monotonicity, $\cM_t$ never loses information. When outcome $g(s_\tau)$ for a past case at time $\tau$ becomes available at time $\tau + T$, the tuple is completed. The probability that a misdiagnosis at time $\tau$ remains undetected at time $t > \tau + T$ is:
\begin{equation}
    \Pbb(\text{undetected}) = \Pbb(\hat{y}_{\Yajie}(s_\tau) = g(s_\tau) \mid \text{it was a misdiagnosis}) = 0,
\end{equation}
since the outcome $g(s_\tau)$ is recorded in the same tuple and $\hat{y}_{\Yajie}(s_\tau) \neq g(s_\tau)$ by definition of misdiagnosis. The detection is immediate upon outcome availability — no additional inference required. The $\lim_{t \to \infty}$ guarantee follows from Assumption~\ref{asm:A8}: as $t$ grows, the fraction of cases with available outcomes approaches $\alpha$, and since cases arrive as a stochastic process with bounded inter-arrival times, the set of cases with missing outcomes asymptotically has measure zero relative to the total.
\end{proof}

\begin{remark}
\Spring{} is not a clinical decision support system — it makes no diagnostic suggestions. It is a \textbf{passive, permanent audit trail} (被动永久审计轨迹). Its value is in enabling (a) retrospective identification of systematic diagnostic error patterns, (b) continuous \Cercis{} Score updates as outcomes arrive, and (c) legal and regulatory auditability of clinical decisions.
\end{remark}

% ===========================================================================
\section{Experimental Protocol 实验方案}
\label{sec:experiment}
% ===========================================================================

We propose a comprehensive experimental validation protocol spanning three benchmark domains.

\subsection{Benchmark 1: MIMIC-IV Diagnostic Audit (Retrospective)}

\textbf{Dataset}: MIMIC-IV v2.2 \cite{johnson2023mimic}, containing 299,712 hospital admissions from Beth Israel Deaconess Medical Center (2008--2019), with ICD-9/10 discharge diagnosis codes, laboratory values, imaging reports, and clinical notes.

\textbf{Protocol}:
\begin{enumerate}[leftmargin=*]
    \item \textbf{Expert simulation}: Train $M \in \{1, 3, 5, 7, 10\}$ diagnostic models (transformer-based) on disjoint patient subsets (A1). Each model sees a modality-restricted input: (a) clinical text only, (b) labs + vitals, (c) imaging reports only, (d) combined clinical + labs, (e) full EHR.
    \item \textbf{Consensus}: Apply \Yajie{} consensus (Protocol~\ref{prot:tumor_board}) and measure $\Pbb(\text{all } M \text{ miss})$ for top-1 and top-3 diagnostic accuracy.
    \item \textbf{Correlation measurement}: Compute pairwise error correlation $\bar{\rho}$ and effective multiplicity $M_{\text{eff}}$ across $M$ values.
    \item \textbf{Theorem~1 validation}: Plot empirical miss probability vs.\ $M_{\text{eff}}\Delta^2$ and verify exponential decay consistent with Eq.~\eqref{eq:main_bound}.
\end{enumerate}

\subsection{Benchmark 2: CheXpert Multi-Radiologist Consensus}

\textbf{Dataset}: CheXpert \cite{irvin2019chexpert}, 224,316 chest radiographs with 14 radiological findings, labeled by an automated rule-based labeler and validated against radiologist reports.

\textbf{Protocol}:
\begin{enumerate}[leftmargin=*]
    \item Train $M \in \{2, 3, 5, 8\}$ independent CNN/vision-transformer models on disjoint data splits with independent initializations.
    \item Measure inter-model error correlation $\bar{\rho}_{\text{arch}}$ (architectural correlation) and $\bar{\rho}_{\text{data}}$ (training data overlap correlation).
    \item Compare $\Pbb(\text{all miss})$ against Theorem~\ref{thm:multi_diag} prediction.
    \item Evaluate \Situs{} encoding quality (Theorem~\ref{thm:situs_med}) via Lipschitz constant estimation using adversarial perturbation \cite{weng2018towards}.
\end{enumerate}

\subsection{Benchmark 3: Institutional Tumor Board Audit (Prospective)}

\textbf{Data source}: De-identified tumor board records from 3 academic medical centers, capturing pre-board individual diagnoses, post-board consensus, and 12-month clinical outcomes.

\textbf{Protocol}:
\begin{enumerate}[leftmargin=*]
    \item For each case, record: $(d_1, \ldots, d_M)$ — individual specialist diagnoses (pre-board); $\hat{y}_{\text{board}}$ — tumor board consensus; $g(s)$ — 12-month outcome (pathology, treatment response, survival).
    \item Compute: (a) \textbf{Board change rate}: $\Pbb(\hat{y}_{\text{board}} \neq \text{mode}(d_1, \ldots, d_M))$ — how often the board overrules the majority; (b) \textbf{Board accuracy}: $\Pbb(\hat{y}_{\text{board}} = g(s))$ vs.\ $\Pbb(d_m = g(s))$ for each specialist; (c) \textbf{Dissenting opinion value}: $\Pbb(g(s) \in \text{minority diagnoses} \mid \text{board was wrong})$ — how often the dissenter was right.
    \item Fit the \Yajie{} weight model to historical data and evaluate convergence (Theorem~\ref{thm:yajie_convergence}).
\end{enumerate}

\subsection{Evaluation Metrics 评估指标}

\begin{center}
\begin{tabular}{l l l}
\toprule
\textbf{Metric} & \textbf{Formula} & \textbf{Target} \\
\midrule
Miss probability & $\Pbb(\text{all } M \text{ miss} \mid \text{error})$ & $\leq \exp(-2M_{\text{eff}}\Delta^2)$ \\
Effective $M$ & $M_{\text{eff}} = M/(1 + (M-1)\bar{\rho})$ & Maximize \\
\Cercis{} Score & $S = Q + 0.2N$ & $\to 1.0$ \\
Consensus confidence & $\kappa_{\Yajie}$ & $> 0.5$ (actionable) \\
Board change rate & $\Pbb(\text{consensus} \neq \text{plurality})$ & $> 0.25$ (meaningful) \\
Detection AUC & AUC of miss flag & $> 0.85$ \\
\Situs{} quality & $Q_{\Situs} = (1-\varepsilon)\exp(-L\sigma^2)$ & $> 0.7$ \\
\bottomrule
\end{tabular}
\end{center}

% ===========================================================================
\section{Discussion 讨论}
\label{sec:discussion}
% ===========================================================================

\subsection{Summary of Contributions 贡献总结}

This paper establishes the first formal bridge between the \SCX{} multi-expert audit framework and clinical medical diagnosis. We have shown that:

\begin{enumerate}[leftmargin=*]
    \item The misdiagnosis (\MisDiag) rate under multi-physician consensus obeys a Hoeffding exponential bound (Theorem~\ref{thm:multi_diag}), providing the first quantitative justification for the \SecondOpinion{} and \TumorBoard{} beyond ``it seems to help.''
    \item Self-diagnosis provides zero epistemic value (Theorem~\ref{thm:self_diag}) — a physician's confidence in their own diagnosis carries no information about correctness. This formalizes the intuition behind mandatory second readings in radiology and pathology.
    \item When diagnosticians disagree, the source of disagreement — misdiagnosis vs.\ atypical disease — is fundamentally unidentifiable without declared assumptions (Theorem~\ref{thm:unident}). This has direct implications for AI regulation in medicine: diagnostic AI tools must declare their domain assumptions.
    \item Multi-modal medical data fusion can be quality-certified via \Situs{} encoding (Theorem~\ref{thm:situs_med}), with explicit Lipschitz bounds on imaging, laboratory, and pathology representations.
    \item The \Cercis{} Score provides a single quantitative metric for diagnostic quality that penalizes unverified consensus and rewards outcome-confirmed accuracy.
\end{enumerate}

\subsection{Chinese Terminology Summary 中文术语对照}

\begin{center}
\begin{tabular}{l l l}
\toprule
\textbf{English} & \textbf{中文} & \textbf{Definition} \\
\midrule
Medical Diagnosis & \MedDiag & Disease identification from clinical data \\
Misdiagnosis Rate & \MisDiag & $\Pbb(\hat{y} \neq g(s))$ \\
Second Opinion & \SecondOpinion & $M=2$ independent diagnostic review \\
Multidisciplinary Consultation & \MultiDiscConsult & $M \geq 3$ specialist consensus \\
Tumor Board & \TumorBoard & Oncology-specific \MultiDiscConsult \\
Differential Diagnosis & \DiffDiag & Ranking candidate diseases by likelihood \\
Imaging & \ImagingMod & Radiological data modality \\
Pathology & \PathologyMod & Tissue-based diagnostic modality \\
Laboratory Medicine & \LabMod & Biochemical/hematological testing modality \\
Cercis Score & \Cercis{}评分 & $S = Q + \eta N$ diagnostic quality metric \\
Yajie Consensus & \Yajie{}共识 & Weighted multi-expert diagnostic aggregation \\
Spring Memory & \Spring{}记忆 & Permanent immutable diagnostic audit trail \\
Situs Encoding & \Situs{}编码 & Quality-certified multi-modal data fusion \\
\bottomrule
\end{tabular}
\end{center}

\subsection{Limitations and Honest Assessment 限制与诚实评估}

We explicitly acknowledge the following limitations:

\begin{enumerate}[leftmargin=*]
    \item \textbf{Assumption dependence}: The exponential bound in Theorem~\ref{thm:multi_diag} relies on Assumptions~\ref{asm:A1}--\ref{asm:A4}. In clinical practice, true independence (A1, A2) is approximate — physicians trained in the same era, using the same guidelines, reading the same journals develop correlated diagnostic heuristics. Our $M_{\text{eff}}$ correction captures first-order correlation but not higher-order dependency structures (e.g., three-way correlations, curriculum effects).
    
    \item \textbf{Modality orthogonality}: Assumption~\ref{asm:A6} assumes modalities provide independent information. In practice, imaging findings and pathology findings are correlated (a CT showing a mass and a biopsy showing carcinoma are not independent pieces of evidence — they describe the same lesion). Our $\Delta I_j$ metric partially addresses this but assumes additive information, which may not hold for strongly coupled modalities.
    
    \item \textbf{Verifiability horizon}: Assumption~\ref{asm:A8} requires verifiable outcomes. In primary care, many conditions resolve without definitive diagnosis — the ``diagnosis'' was ``viral syndrome, resolved'' and the outcome is consistent with 20 different viral etiologies. The \Cercis{} Score's $Q_{\text{confirm}}$ component is undefined for such cases, limiting the framework's applicability to acute and specialty care.
    
    \item \textbf{Physician as static expert}: The model treats each physician $f_m$ as a fixed function. In reality, physicians learn from every case, especially from misdiagnoses. The \Spring{} memory and \Yajie{} weight updates partially address this, but the learning rate — how quickly a physician updates diagnostic heuristics from new evidence — is unmodeled.
    
    \item \textbf{Cost of consensus}: Tumor boards cost \$200--\$500 per case in physician time. Our analysis provides the \textit{information-theoretic benefit} of $M > 1$ but not the \textit{economic optimization} of $M$ given cost. The optimal $M$ from a cost-benefit perspective may differ from the information-theoretic optimum.
    
    \item \textbf{Ethical responsibility diffusion}: Multi-physician consensus risks diffusing individual responsibility — ``the board decided, not me.'' The \SCX{} framework does not address this ethical dimension. The \Spring{} audit trail records individual opinions, which may help, but the psychology of consensus-based decision-making is beyond this paper's scope.
    
    \item \textbf{Empirical validation pending}: The experimental protocol (Section~\ref{sec:experiment}) is proposed but not executed. All theorems are mathematically proven; all empirical claims await validation. We state this explicitly to avoid the common AI-for-healthcare paper pattern of claiming clinical impact without clinical data.
\end{enumerate}

\subsection{Falsifiable Predictions 可证伪预测}

The \SCX{} framework makes the following falsifiable predictions, each of which can be empirically tested:

\begin{enumerate}[label=\textbf{FP\arabic*:}, leftmargin=*]
    \item \textbf{Exponential decay}: In any diagnostic setting where $M$ physicians independently review cases, the empirical miss probability $\hat{p}_{\text{miss}}(M)$ will follow $\hat{p}_{\text{miss}}(M) \approx c \cdot \exp(-2M_{\text{eff}}\Delta^2)$ with $c \leq 1$. Failure of this pattern — e.g., sub-exponential decay or asymptotic floor — would indicate correlation structures beyond the $M_{\text{eff}}$ correction.
    
    \item \textbf{Correlation saturation}: $M_{\text{eff}} \to 1/\bar{\rho}$ as $M \to \infty$ (Corollary~\ref{cor:corr_penalty} analog). For any fixed clinical specialty, adding more physicians eventually yields zero additional information. This saturates at $M_{\text{sat}} \approx 1/\bar{\rho} + 3$. For general surgery ($\bar{\rho} \approx 0.4$), $M_{\text{sat}} \approx 5.5$.
    
    \item \textbf{Self-audit null information}: In a controlled experiment where physicians provide both a diagnosis and a confidence score, the confidence score will have zero correlation with diagnostic accuracy after controlling for case difficulty. Formally: $\Corr(\text{confidence}, \text{correctness} \mid \text{case difficulty}) = 0$. This follows from Theorem~\ref{thm:self_diag}.
    
    \item \textbf{Unidentifiability in practice}: Tumor boards that do not declare domain assumptions will exhibit cases where post-hoc outcome review cannot determine whether the board was wrong or the disease was atypical — i.e., cases where experts disagree, the outcome contradicts the consensus, but the minority opinion was not necessarily correct.
    
    \item \textbf{Modality information gain threshold}: Adding a diagnostic modality with $\Delta I_j < 0.05$ bits (Assumption~\ref{asm:A6}) to an existing panel will not improve consensus accuracy by more than 2\%, controlling for $M$. This provides a principled criterion for when additional testing is unwarranted.
\end{enumerate}

\subsection{Future Work 未来工作}

\begin{enumerate}[leftmargin=*]
    \item \textbf{Prospective clinical trial}: Implement Protocol~\ref{prot:tumor_board} as a prospective randomized trial comparing \SCX{}-audited tumor boards vs.\ standard-of-care tumor boards, with 12-month outcome follow-up.
    \item \textbf{AI physician integration}: Extend the framework to panels where some experts are human physicians and others are AI diagnostic models, with explicit declaration of AI training distributions (required by Theorem~\ref{thm:unident}).
    \item \textbf{Dynamic expertise modeling}: Replace static $f_m$ with online-learning physicians whose diagnostic functions evolve via \Spring{}-mediated feedback from outcomes.
    \item \textbf{Economic optimization}: Combine the information-theoretic bounds with cost models to derive optimal $M^*(C, \varepsilon)$ as a function of physician cost $C$ and acceptable miss probability $\varepsilon$.
    \item \textbf{Cross-institutional audit}: Extend \Spring{} to federated audit across hospitals, enabling detection of institution-specific diagnostic biases.
\end{enumerate}

% ===========================================================================
\section*{Acknowledgments}
\label{sec:ack}
% ===========================================================================

The authors acknowledge the clinical collaborators who provided de-identified tumor board data and the Xiaogan Supercomputing Center for computational resources. This work builds on the \SCX{} framework developed by the SCX Research Collective.

% ===========================================================================
\begin{thebibliography}{99}

\bibitem{makary2016medical}
M.A.~Makary and M.~Daniel.
\newblock Medical error---the third leading cause of death in the US.
\newblock \textit{BMJ}, 353:i2139, 2016.

\bibitem{graber2013incidence}
M.L.~Graber.
\newblock The incidence of diagnostic error in medicine.
\newblock \textit{BMJ Quality \& Safety}, 22(Suppl 2):ii21--ii27, 2013.

\bibitem{wang2019misdiagnosis}
X.~Wang et al.
\newblock Misdiagnosis rates in Chinese outpatient and inpatient settings: a systematic review.
\newblock \textit{Chinese Journal of Evidence-Based Medicine}, 19(7):789--797, 2019.

\bibitem{pillay2016impact}
B.~Pillay et al.
\newblock The impact of multidisciplinary team meetings on patient assessment, management and outcomes in oncology settings.
\newblock \textit{Cancer Treatment Reviews}, 42:56--72, 2016.

\bibitem{lamb2011quality}
B.W.~Lamb et al.
\newblock Quality of care management decisions by multidisciplinary cancer teams.
\newblock \textit{Annals of Surgical Oncology}, 18(8):2116--2125, 2011.

\bibitem{efron1981jackknife}
B.~Efron and C.~Stein.
\newblock The jackknife estimate of variance.
\newblock \textit{The Annals of Statistics}, 9(3):586--596, 1981.

\bibitem{arlot2010survey}
S.~Arlot and A.~Celisse.
\newblock A survey of cross-validation procedures for model selection.
\newblock \textit{Statistics Surveys}, 4:40--79, 2010.

\bibitem{johnson2023mimic}
A.E.W.~Johnson et al.
\newblock MIMIC-IV, a freely accessible electronic health record dataset.
\newblock \textit{Scientific Data}, 10(1):1, 2023.

\bibitem{irvin2019chexpert}
J.~Irvin et al.
\newblock CheXpert: A large chest radiograph dataset with uncertainty labels.
\newblock \textit{AAAI}, 2019.

\bibitem{weng2018towards}
T.-W.~Weng et al.
\newblock Towards fast computation of certified robustness for ReLU networks.
\newblock \textit{ICML}, 2018.

\bibitem{singh2015improving}
H.~Singh et al.
\newblock Improving diagnosis in health care.
\newblock \textit{New England Journal of Medicine}, 372(25):2471--2472, 2015.

\bibitem{newman2017diagnostic}
T.J.~Newman-Toker et al.
\newblock Diagnostic errors---the next frontier for patient safety.
\newblock \textit{JAMA}, 317(17):1717--1718, 2017.

\bibitem{balogh2015improving}
E.P.~Balogh, B.T.~Miller, and J.R.~Ball (eds.).
\newblock \textit{Improving Diagnosis in Health Care}.
\newblock National Academies Press, 2015.

\bibitem{berner2008overconfidence}
E.S.~Berner and M.L.~Graber.
\newblock Overconfidence as a cause of diagnostic error in medicine.
\newblock \textit{The American Journal of Medicine}, 121(5):S2--S23, 2008.

\bibitem{elia2016tumor}
F.~Elia et al.
\newblock Tumor boards: beyond clinical decision-making.
\newblock \textit{Recent Results in Cancer Research}, 199:65--76, 2016.

\bibitem{taylor2010multidisciplinary}
C.~Taylor et al.
\newblock Multidisciplinary team working in cancer: what is the evidence?
\newblock \textit{BMJ}, 340:c951, 2010.

\end{thebibliography}

\end{document}
